\section{Document de Conception
Technique}\label{document-de-conception-technique}

\subsection{\texorpdfstring{\texttt{Inventaires\_completude\_representativite.R}}{Inventaires\_completude\_representativite.R}}\label{inventaires_completude_representativite.r}

\begin{center}\rule{0.5\linewidth}{0.5pt}\end{center}

\textbf{Projet} : Projet hôte / Inventaires fongiques\\
\textbf{Auteur} : Eddy Boite\\
\textbf{Version script} : 1.4 (selon en-tête du script)\\
\textbf{Date doc} : Mars 2026\\
\textbf{Type de document} : Spécification technique

\begin{center}\rule{0.5\linewidth}{0.5pt}\end{center}

\subsection{Table des matières}\label{table-des-matiuxe8res}

\begin{enumerate}
\def\labelenumi{\arabic{enumi}.}
\tightlist
\item
  \hyperref[1-vue-densemble]{Vue d'ensemble}
\item
  \hyperref[2-architecture-du-systux5cux25C3ux5cux25A8me]{Architecture
  du système}
\item
  \hyperref[3-analyses-de-complux5cux25C3ux5cux25A9tude]{Analyses de
  complétude}
\item
  \hyperref[4-analyses-de-reprux5cux25C3ux5cux25A9sentativitux5cux25C3ux5cux25A9]{Analyses
  de représentativité}
\item
  \hyperref[5-configuration-rux5cux25C3ux5cux25A9elle-du-script]{Configuration
  réelle du script}
\item
  \hyperref[6-portabilitux5cux25C3ux5cux25A9-et-nom-du-projet]{Portabilité
  et nom du projet}
\item
  \hyperref[7-journalisation-et-traux5cux25C3ux5cux25A7abilitux5cux25C3ux5cux25A9]{Journalisation
  et traçabilité}
\item
  \hyperref[8-validation-csv-et-conformitux5cux25C3ux5cux25A9]{Validation
  CSV et conformité}
\item
  \hyperref[9-rapports-et-livrables]{Rapports et livrables}
\item
  \hyperref[9-flux-de-traitement]{Flux de traitement}
\item
  \hyperref[10-structures-de-donnux5cux25C3ux5cux25A9es-exhaustif]{Structures
  de données (exhaustif)}
\item
  \hyperref[11-mux5cux25C3ux5cux25A9triques-et-interprux5cux25C3ux5cux25A9tation]{Métriques
  et interprétation}
\item
  \hyperref[12-versioning-et-maintenance]{Versioning et maintenance}
\item
  \hyperref[13-exux5cux25C3ux5cux25A9cution-dux5cux25C3ux5cux25A9pendances-et-limites]{Exécution,
  dépendances et limites}
\item
  \hyperref[14-annexes]{Annexes}
\end{enumerate}

\begin{center}\rule{0.5\linewidth}{0.5pt}\end{center}

\subsection{1. Vue d'ensemble}\label{vue-densemble}

\subsubsection{1.1 Objectif}\label{objectif}

Le script \texttt{scripts/Inventaires\_completude\_representativite.R}
automatise l'analyse de complétude et de représentativité des
inventaires fongiques, par site et globalement.

Il produit notamment :

\begin{itemize}
\tightlist
\item
  la courbe temps-espèces,
\item
  l'ajustement linéaire/hyperbolique (si possible),
\item
  l'indice \texttt{TEE} (taux d'espèces exceptionnelles) et
  \texttt{Ir\ =\ 1\ -\ TEE},
\item
  les fréquences d'occurrence par espèce,
\item
  les métriques de pertinence scientifique,
\item
  les comparaisons inter-sites.
\end{itemize}

\subsubsection{1.2 Données manipulées (clarification
importante)}\label{donnuxe9es-manipuluxe9es-clarification-importante}

Le script utilise la colonne taxonomique \texttt{espece}.

\begin{itemize}
\tightlist
\item
  Clés d'analyse principales : \texttt{site}, \texttt{date},
  \texttt{visite\_id}, \texttt{espece}.
\item
  Clé de visite : \texttt{date\ +\ visite\_id} (et \texttt{site} en
  multi-sites).
\end{itemize}

\subsubsection{1.3 Entrées / sorties de haut
niveau}\label{entruxe9es-sorties-de-haut-niveau}

\textbf{Entrée principale} :

\begin{itemize}
\tightlist
\item
  \texttt{data/observations.csv} (ou \texttt{observations.tsv},
  \texttt{observations.txt})
\end{itemize}

\textbf{Sorties globales} :

\begin{itemize}
\tightlist
\item
  \texttt{results/ICR\_donnees\_preparees.csv}
\item
  \texttt{results/ICR\_resume\_tous\_sites.csv}
\item
  \texttt{results/ICR\_metrics\_pertinence\_tous\_sites.csv}
\item
  graphiques globaux \texttt{results/ICR\_*.png}
\end{itemize}

\textbf{Sorties par site}
(\texttt{results/\textless{}site\textgreater{}/}) :

\begin{itemize}
\tightlist
\item
  \texttt{ICR\_01} à \texttt{ICR\_07} en CSV et PNG selon disponibilité
  des données/modèles.
\end{itemize}

\begin{center}\rule{0.5\linewidth}{0.5pt}\end{center}

\subsection{2. Architecture du
système}\label{architecture-du-systuxe8me}

\subsubsection{2.1 Modules fonctionnels}\label{modules-fonctionnels}

Le script suit la chaîne suivante :

\begin{enumerate}
\def\labelenumi{\arabic{enumi}.}
\tightlist
\item
  Initialisation (options, packages, constantes)
\item
  Validation configuration
\item
  Résolution du fichier d'entrée
\item
  Lecture auto-délimiteur
\item
  Préparation/validation des données (\texttt{prepare\_data})
\item
  Analyse par site (\texttt{analyze\_site})
\item
  Consolidation globale (\texttt{run\_analysis})
\item
  Export CSV/PNG + manifeste de sorties
\end{enumerate}

\subsubsection{2.2 Fonctions clés}\label{fonctions-cluxe9s}

\begin{itemize}
\tightlist
\item
  \texttt{read\_delim\_auto()} : détection automatique du séparateur
  (\texttt{,} \texttt{;} \texttt{TAB}) + récupération des problèmes de
  parsing.
\item
  \texttt{build\_csv\_colspec()} : schéma de lecture CSV (colonnes
  attendues en \texttt{character}).
\item
  \texttt{audit\_csv\_conformity()} : audit de conformité (colonnes
  requises, colonnes extra, parsing).
\item
  \texttt{export\_csv\_conformity\_report()} : export des rapports
  \texttt{ICR\_00\_*}.
\item
  \texttt{prepare\_data()} : validation stricte + déduplication.
\item
  \texttt{calc\_cumulative\_metrics()} : richesse par visite et cumul.
\item
  \texttt{calc\_tee\_ir()} : calcul dynamique de
  \texttt{TEE}/\texttt{Ir}.
\item
  \texttt{calc\_species\_frequency()} : fréquence temporelle par espèce.
\item
  \texttt{calc\_spatial\_occupancy()} : occupation spatiale (si
  \texttt{placette}).
\item
  \texttt{fit\_linear\_model()} / \texttt{fit\_hyperbolic\_model()} :
  modélisation de la courbe.
\item
  \texttt{calc\_scientific\_metrics()} : score synthétique et métriques
  de robustesse.
\item
  \texttt{log\_output\_manifest()} : contrôle d'intégrité des livrables.
\end{itemize}

\begin{center}\rule{0.5\linewidth}{0.5pt}\end{center}

\subsection{3. Analyses de complétude}\label{analyses-de-compluxe9tude}

\subsubsection{3.1 Richesse observée et
cumulée}\label{richesse-observuxe9e-et-cumuluxe9e}

Par visite, le script calcule :

\begin{itemize}
\tightlist
\item
  \texttt{nb\_especes\_trouvees}
\item
  \texttt{nb\_nouvelles}
\item
  \texttt{nb\_cumule}
\end{itemize}

Sorties :

\begin{itemize}
\tightlist
\item
  \texttt{ICR\_01\_courbe\_temps\_especes.csv}
\item
  \texttt{ICR\_01\_richesse\_par\_visite\_et\_cumul.png}
\end{itemize}

\subsubsection{3.2 Complétude observée /
asymptote}\label{compluxe9tude-observuxe9e-asymptote}

Quand l'ajustement hyperbolique est disponible :

\textbf{Principe} : la complétude correspond au rapport entre la
richesse observée et l'asymptote estimée.

avec \texttt{S\_asymptote\ =\ a} (paramètre du modèle hyperbolique).

Conditions d'estimation :

\begin{itemize}
\tightlist
\item
  \texttt{nb\_visites\ \textgreater{}=\ CONFIG\$min\_visits\_for\_model}
\item
  package optionnel \texttt{minpack.lm} installé.
\end{itemize}

Sorties :

\begin{itemize}
\tightlist
\item
  \texttt{ICR\_05\_modeles.csv} (si modèle ajusté)
\item
  \texttt{ICR\_04\_resume\_site.csv}
\item
  \texttt{ICR\_resume\_tous\_sites.csv}
\end{itemize}

\subsubsection{3.3 Fréquence d'occurrence des
espèces}\label{fruxe9quence-doccurrence-des-espuxe8ces}

Le script estime :

\begin{itemize}
\tightlist
\item
  \texttt{nb\_visites} par espèce,
\item
  \texttt{prop\_visites\ =\ nb\_visites\ /\ nb\_visites\_total},
\item
  classe de fréquence via \texttt{CONFIG\$freq\_breaks} +
  \texttt{CONFIG\$freq\_labels}.
\end{itemize}

Sorties :

\begin{itemize}
\tightlist
\item
  \texttt{ICR\_03\_frequence\_especes.csv}
\item
  \texttt{ICR\_04\_histogramme\_frequences.png}
\end{itemize}

\subsubsection{3.4 Couverture spatiale
(optionnelle)}\label{couverture-spatiale-optionnelle}

Si \texttt{placette} est exploitable :

\begin{itemize}
\tightlist
\item
  \texttt{nb\_placettes} par espèce,
\item
  \texttt{prop\_placettes}.
\end{itemize}

Sorties :

\begin{itemize}
\tightlist
\item
  \texttt{ICR\_06\_occupation\_spatiale.csv}
\item
  \texttt{ICR\_05\_temporel\_vs\_spatial.png}
\end{itemize}

\begin{center}\rule{0.5\linewidth}{0.5pt}\end{center}

\subsection{4. Analyses de
représentativité}\label{analyses-de-repruxe9sentativituxe9}

\subsubsection{4.1 TEE et indice Ir}\label{tee-et-indice-ir}

À chaque visite :

\begin{itemize}
\tightlist
\item
  \texttt{TEE} = proportion d'espèces observées une seule fois,
\item
  \texttt{Ir\ =\ 1\ -\ TEE}.
\end{itemize}

Sorties :

\begin{itemize}
\tightlist
\item
  \texttt{ICR\_02\_tee\_ir.csv}
\item
  \texttt{ICR\_03\_tee\_ir.png}
\item
  \texttt{ICR\_comparaison\_ir\_sites.png}
\end{itemize}

\subsubsection{4.2 Stabilité terminale}\label{stabilituxe9-terminale}

Le script calcule :

\begin{itemize}
\tightlist
\item
  \texttt{slope\_terminal\_cumul} (pente des dernières visites),
\item
  \texttt{tail\_novelty\_share} (part des nouveautés en fin de série).
\end{itemize}

Objectif : quantifier la fermeture de l'inventaire.

\subsubsection{4.3 Cohérence
temporelle/spatiale}\label{cohuxe9rence-temporellespatiale}

Si \texttt{placette} existe :

\begin{itemize}
\tightlist
\item
  corrélation de Spearman \texttt{spearman\_temporal\_spatial} entre
  \texttt{prop\_visites} et \texttt{prop\_placettes},
\item
  \texttt{discordance\_rate}.
\end{itemize}

\subsubsection{4.4 Score de pertinence
scientifique}\label{score-de-pertinence-scientifique}

Le score agrégé \texttt{score\_pertinence} utilise les pondérations
codées dans le script :

\begin{itemize}
\tightlist
\item
  \texttt{Ir\ final} : 35 \%
\item
  \texttt{complétude} : 35 \%
\item
  \texttt{stabilité\ finale} : 20 \%
\item
  \texttt{robustesse\ modèle} : 10 \%
\end{itemize}

Sorties :

\begin{itemize}
\tightlist
\item
  \texttt{ICR\_07\_metrics\_pertinence.csv}
\item
  \texttt{ICR\_metrics\_pertinence\_tous\_sites.csv}
\item
  \texttt{ICR\_07\_metrics\_pertinence\_dashboard.png}
\item
  \texttt{ICR\_metrics\_pertinence\_heatmap\_sites.png}
\item
  \texttt{ICR\_metrics\_pertinence\_score\_sites.png}
\end{itemize}

\begin{center}\rule{0.5\linewidth}{0.5pt}\end{center}

\subsection{5. Configuration réelle du
script}\label{configuration-ruxe9elle-du-script}

\subsubsection{5.1 Paramètres runtime (dans le
code)}\label{paramuxe8tres-runtime-dans-le-code}

Le script n'utilise pas de \texttt{.Renviron} dédié ; la configuration
est portée par la liste \texttt{CONFIG}.

{\def\LTcaptype{none} % do not increment counter
\begin{longtable}[]{@{}
  >{\raggedright\arraybackslash}p{(\linewidth - 4\tabcolsep) * \real{0.3333}}
  >{\raggedright\arraybackslash}p{(\linewidth - 4\tabcolsep) * \real{0.3333}}
  >{\raggedright\arraybackslash}p{(\linewidth - 4\tabcolsep) * \real{0.3333}}@{}}
\toprule\noalign{}
\begin{minipage}[b]{\linewidth}\raggedright
Clé CONFIG
\end{minipage} & \begin{minipage}[b]{\linewidth}\raggedright
Défaut
\end{minipage} & \begin{minipage}[b]{\linewidth}\raggedright
Rôle
\end{minipage} \\
\midrule\noalign{}
\endhead
\bottomrule\noalign{}
\endlastfoot
\texttt{input\_file} & \texttt{data/observations.csv} & Fichier
d'entrée \\
\texttt{output\_dir} & \texttt{results/ICR} & Répertoire de sortie \\
\texttt{date\_format} & \texttt{\%Y-\%m-\%d} & Format de parsing des
dates \\
\texttt{csv\_strict\_mode} & \texttt{TRUE} & Stoppe l'exécution si audit
CSV non conforme \\
\texttt{csv\_allow\_extra\_cols} & \texttt{FALSE} & Autorise ou non les
colonnes supplémentaires \\
\texttt{csv\_required\_cols} &
\texttt{c("site","date","visite\_id","espece")} & Colonnes obligatoires
attendues \\
\texttt{csv\_optional\_cols} & \texttt{c("placette")} & Colonnes
optionnelles autorisées \\
\texttt{freq\_breaks} & \texttt{c(0,0.10,0.25,0.50,0.75,1.00)} & Bornes
de classes de fréquence \\
\texttt{freq\_labels} &
\texttt{exceptionnelle}\ldots{}\texttt{constante} & Labels des
classes \\
\texttt{make\_ca} & \texttt{TRUE} & Active CA/AFC (si \texttt{vegan}) \\
\texttt{min\_visits\_for\_model} & \texttt{5} & Seuil mini pour
modélisation \\
\texttt{width},\texttt{height},\texttt{dpi} &
\texttt{10},\texttt{6},\texttt{300} & Paramètres d'export graphique \\
\end{longtable}
}

\subsubsection{5.2 Flags d'exécution}\label{flags-dexuxe9cution}

{\def\LTcaptype{none} % do not increment counter
\begin{longtable}[]{@{}
  >{\raggedright\arraybackslash}p{(\linewidth - 4\tabcolsep) * \real{0.3333}}
  >{\raggedright\arraybackslash}p{(\linewidth - 4\tabcolsep) * \real{0.3333}}
  >{\raggedright\arraybackslash}p{(\linewidth - 4\tabcolsep) * \real{0.3333}}@{}}
\toprule\noalign{}
\begin{minipage}[b]{\linewidth}\raggedright
Variable/flag
\end{minipage} & \begin{minipage}[b]{\linewidth}\raggedright
Défaut
\end{minipage} & \begin{minipage}[b]{\linewidth}\raggedright
Description
\end{minipage} \\
\midrule\noalign{}
\endhead
\bottomrule\noalign{}
\endlastfoot
\texttt{DEBUG\_MODE} & \texttt{FALSE} & Verbosité et affichage des
avertissements \\
\texttt{BENCHMARK\_MODE} & \texttt{FALSE} & Chronométrage des calculs \\
\texttt{CLEAN\_ENVIRONMENT} & \texttt{FALSE} & Nettoyage environnement
(usage expert) \\
\texttt{INVENTAIRES\_AUTO\_RUN} (env) & \texttt{TRUE} & Exécute
automatiquement \texttt{run\_analysis()} \\
\end{longtable}
}

\begin{center}\rule{0.5\linewidth}{0.5pt}\end{center}

\subsection{6. Portabilité et nom du
projet}\label{portabilituxe9-et-nom-du-projet}

\subsubsection{6.1 Principe}\label{principe}

Le script est \textbf{indépendant du nom du dépôt} : aucun identifiant
de projet (ex. \texttt{statistics}) n'est requis dans le code d'analyse.

La portabilité repose sur :

\begin{itemize}
\tightlist
\item
  des chemins relatifs (\texttt{data/}, \texttt{scripts/},
  \texttt{results/}),
\item
  la résolution dynamique du dossier projet (\texttt{PROJECT\_DIR}),
\item
  des variables d'environnement pour adapter les chemins/fomat de date.
\end{itemize}

\subsubsection{6.2 Variables d'environnement à
utiliser}\label{variables-denvironnement-uxe0-utiliser}

Le nom du projet se configure indirectement en adaptant les chemins à
votre repo :

\begin{itemize}
\tightlist
\item
  \texttt{INVENTAIRES\_INPUT\_FILE} (ex.
  \texttt{data/observations\_public.csv})
\item
  \texttt{INVENTAIRES\_OUTPUT\_DIR} (ex.
  \texttt{results/public\_release})
\item
  \texttt{INVENTAIRES\_DATE\_FORMAT} (ex. \texttt{\%d/\%m/\%Y} si
  nécessaire)
\item
  \texttt{INVENTAIRES\_CSV\_STRICT} (mode strict conformité CSV)
\item
  \texttt{INVENTAIRES\_CSV\_ALLOW\_EXTRA\_COLS} (tolérance des colonnes
  additionnelles)
\end{itemize}

Ainsi, un même script peut être déplacé d'un projet à l'autre sans
modification du code source.

\subsubsection{6.3 Recommandation de
documentation}\label{recommandation-de-documentation}

Pour la diffusion, employer systématiquement un placeholder explicite :

\begin{itemize}
\tightlist
\item
  \texttt{\textless{}nom\_du\_projet\textgreater{}} pour l'arborescence,
\item
  \texttt{\textless{}nom\_du\_projet\textgreater{}.Rproj} pour le
  fichier RStudio,
\item
  chemins d'exécution relatifs (\texttt{Rscript\ scripts/...}).
\end{itemize}

\begin{center}\rule{0.5\linewidth}{0.5pt}\end{center}

\subsection{7. Journalisation et
traçabilité}\label{journalisation-et-trauxe7abilituxe9}

\subsubsection{6.1 Mécanisme réel}\label{muxe9canisme-ruxe9el}

La journalisation est effectuée via les fonctions :

\begin{itemize}
\tightlist
\item
  \texttt{log\_info()}
\item
  \texttt{log\_debug()}
\item
  \texttt{log\_warning()}
\end{itemize}

La sortie se fait sur la \textbf{console standard} ; ce script
n'implémente pas de système de fichiers de logs dédié.

\subsubsection{6.2 Manifeste de sorties}\label{manifeste-de-sorties}

En fin d'exécution, \texttt{log\_output\_manifest()} publie un état des
livrables :

\begin{itemize}
\tightlist
\item
  \texttt{OK}
\item
  \texttt{OPTIONNEL\_NON\_GENERE}
\item
  \texttt{MANQUANT}
\end{itemize}

Ce manifeste sert de contrôle qualité de production.

\begin{center}\rule{0.5\linewidth}{0.5pt}\end{center}

\subsection{8. Validation CSV et
conformité}\label{validation-csv-et-conformituxe9}

\subsubsection{8.1 Contrôles appliqués}\label{contruxf4les-appliquuxe9s}

Avant \texttt{prepare\_data()}, le script exécute un audit de conformité
:

\begin{enumerate}
\def\labelenumi{\arabic{enumi}.}
\tightlist
\item
  vérification des colonnes obligatoires,
\item
  détection de colonnes supplémentaires (selon
  \texttt{csv\_allow\_extra\_cols}),
\item
  comptage des problèmes de parsing \texttt{readr::problems()}.
\end{enumerate}

\subsubsection{8.2 Mode strict}\label{mode-strict}

Si \texttt{csv\_strict\_mode\ =\ TRUE} (défaut), l'exécution est
interrompue en cas de non-conformité.

\subsubsection{8.3 Rapports d'audit}\label{rapports-daudit}

\begin{itemize}
\tightlist
\item
  \texttt{results/ICR/ICR\_00\_csv\_conformite\_report.csv} (toujours)
\item
  \texttt{results/ICR/ICR\_00\_csv\_conformite\_problems.csv} (si
  anomalies parsing)
\end{itemize}

\subsection{9. Rapports et livrables}\label{rapports-et-livrables}

\subsubsection{7.1 CSV globaux (toujours)}\label{csv-globaux-toujours}

\begin{itemize}
\tightlist
\item
  \texttt{results/ICR/ICR\_00\_csv\_conformite\_report.csv}
\item
  \texttt{results/ICR\_donnees\_preparees.csv}
\item
  \texttt{results/ICR\_resume\_tous\_sites.csv}
\item
  \texttt{results/ICR\_metrics\_pertinence\_tous\_sites.csv}
\end{itemize}

\subsubsection{7.1 bis CSV globaux
(conditionnels)}\label{bis-csv-globaux-conditionnels}

\begin{itemize}
\tightlist
\item
  \texttt{results/ICR/ICR\_00\_csv\_conformite\_problems.csv} (si
  parsing anormal)
\end{itemize}

\subsubsection{\texorpdfstring{7.2 CSV par site
(\texttt{results/\textless{}site\textgreater{}/})}{7.2 CSV par site (results/\textless site\textgreater/)}}\label{csv-par-site-resultssite}

\begin{itemize}
\tightlist
\item
  \texttt{ICR\_01\_courbe\_temps\_especes.csv}
\item
  \texttt{ICR\_02\_tee\_ir.csv}
\item
  \texttt{ICR\_03\_frequence\_especes.csv}
\item
  \texttt{ICR\_04\_resume\_site.csv}
\item
  \texttt{ICR\_05\_modeles.csv} (si modèle)
\item
  \texttt{ICR\_06\_occupation\_spatiale.csv} (si \texttt{placette})
\item
  \texttt{ICR\_07\_metrics\_pertinence.csv}
\end{itemize}

\subsubsection{7.3 Graphiques globaux}\label{graphiques-globaux}

\begin{itemize}
\tightlist
\item
  \texttt{ICR\_comparaison\_completude\_sites.png}
\item
  \texttt{ICR\_comparaison\_ir\_sites.png}
\item
  \texttt{ICR\_metrics\_pertinence\_heatmap\_sites.png}
\item
  \texttt{ICR\_metrics\_pertinence\_score\_sites.png}
\end{itemize}

\subsubsection{7.4 Graphiques par site}\label{graphiques-par-site}

\begin{itemize}
\tightlist
\item
  \texttt{ICR\_01\_richesse\_par\_visite\_et\_cumul.png}
\item
  \texttt{ICR\_02\_courbe\_temps\_especes\_hyperbole.png}
\item
  \texttt{ICR\_03\_tee\_ir.png}
\item
  \texttt{ICR\_04\_histogramme\_frequences.png}
\item
  \texttt{ICR\_05\_temporel\_vs\_spatial.png} (si \texttt{placette})
\item
  \texttt{ICR\_06\_CA\_placettes\_especes.png} (si \texttt{make\_ca} +
  \texttt{vegan})
\item
  \texttt{ICR\_07\_metrics\_pertinence\_dashboard.png}
\end{itemize}

\begin{center}\rule{0.5\linewidth}{0.5pt}\end{center}

\subsection{9. Flux de traitement}\label{flux-de-traitement}

\begin{verbatim}
run_analysis(CONFIG)
  ├─ validate_config()
  ├─ resolve_input_file()
  ├─ read_delim_auto()
  ├─ audit_csv_conformity()
  ├─ export_csv_conformity_report()
  ├─ stop si non conforme en mode strict
  ├─ prepare_data()
  ├─ write ICR_donnees_preparees.csv
  ├─ for each site:
  │    └─ analyze_site()
  │         ├─ calc_cumulative_metrics()
  │         ├─ calc_tee_ir()
  │         ├─ calc_species_frequency()
  │         ├─ calc_spatial_occupancy() [optionnel]
  │         ├─ fit models [conditionnel]
  │         ├─ calc_scientific_metrics()
  │         ├─ export CSV site
  │         └─ export PNG site
  ├─ consolidation multi-sites (CSV globaux)
  ├─ export graphiques globaux
  └─ log_output_manifest()
\end{verbatim}

\begin{center}\rule{0.5\linewidth}{0.5pt}\end{center}

\subsection{10. Structures de données
(exhaustif)}\label{structures-de-donnuxe9es-exhaustif}

\subsubsection{9.1 Entrée}\label{entruxe9e}

{\def\LTcaptype{none} % do not increment counter
\begin{longtable}[]{@{}
  >{\raggedright\arraybackslash}p{(\linewidth - 4\tabcolsep) * \real{0.3333}}
  >{\raggedright\arraybackslash}p{(\linewidth - 4\tabcolsep) * \real{0.3333}}
  >{\raggedright\arraybackslash}p{(\linewidth - 4\tabcolsep) * \real{0.3333}}@{}}
\toprule\noalign{}
\begin{minipage}[b]{\linewidth}\raggedright
Fichier
\end{minipage} & \begin{minipage}[b]{\linewidth}\raggedright
Colonnes obligatoires
\end{minipage} & \begin{minipage}[b]{\linewidth}\raggedright
Colonnes optionnelles
\end{minipage} \\
\midrule\noalign{}
\endhead
\bottomrule\noalign{}
\endlastfoot
\texttt{data/observations.csv} (ou \texttt{.tsv}, \texttt{.txt}) &
\texttt{site}, \texttt{date}, \texttt{visite\_id}, \texttt{espece} &
\texttt{placette} \\
\end{longtable}
}

\subsubsection{9.2 Règles de validation
d'entrée}\label{ruxe8gles-de-validation-dentruxe9e}

\begin{enumerate}
\def\labelenumi{\arabic{enumi}.}
\tightlist
\item
  Colonnes obligatoires présentes.
\item
  Colonnes supplémentaires refusées par défaut (configurable).
\item
  Audit des problèmes de parsing (\texttt{readr::problems()}).
\item
  Parsing date selon \texttt{CONFIG\$date\_format}.
\item
  \texttt{site}, \texttt{visite\_id}, \texttt{espece} non vides.
\item
  Déduplication sur (\texttt{site}, \texttt{placette}, \texttt{date},
  \texttt{visite\_id}, \texttt{espece}).
\end{enumerate}

\subsubsection{9.3 Colonnes de sortie
principales}\label{colonnes-de-sortie-principales}

\paragraph{\texorpdfstring{\texttt{ICR\_01\_courbe\_temps\_especes.csv}}{ICR\_01\_courbe\_temps\_especes.csv}}\label{icr_01_courbe_temps_especes.csv}

\begin{itemize}
\tightlist
\item
  \texttt{visite\_index}, \texttt{visite\_id}, \texttt{date},
  \texttt{nb\_especes\_trouvees}, \texttt{nb\_nouvelles},
  \texttt{nb\_cumule}
\end{itemize}

\paragraph{\texorpdfstring{\texttt{ICR\_02\_tee\_ir.csv}}{ICR\_02\_tee\_ir.csv}}\label{icr_02_tee_ir.csv}

\begin{itemize}
\tightlist
\item
  \texttt{visite\_index}, \texttt{visite\_id}, \texttt{date},
  \texttt{nb\_total}, \texttt{nb\_esp\_1fois}, \texttt{TEE}, \texttt{Ir}
\end{itemize}

\paragraph{\texorpdfstring{\texttt{ICR\_03\_frequence\_especes.csv}}{ICR\_03\_frequence\_especes.csv}}\label{icr_03_frequence_especes.csv}

\begin{itemize}
\tightlist
\item
  \texttt{espece}, \texttt{nb\_visites}, \texttt{prop\_visites},
  \texttt{classe\_frequence}
\end{itemize}

\paragraph{\texorpdfstring{\texttt{ICR\_04\_resume\_site.csv}}{ICR\_04\_resume\_site.csv}}\label{icr_04_resume_site.csv}

\begin{itemize}
\tightlist
\item
  \texttt{site}, \texttt{nb\_observations}, \texttt{nb\_visites},
  \texttt{nb\_especes\_observees}, \texttt{nb\_placettes},
  \texttt{asymptote\_hyperbolique},
  \texttt{completude\_obs\_sur\_asymptote}, \texttt{tee\_final},
  \texttt{ir\_final}
\end{itemize}

\paragraph{\texorpdfstring{\texttt{ICR\_05\_modeles.csv}
(conditionnel)}{ICR\_05\_modeles.csv (conditionnel)}}\label{icr_05_modeles.csv-conditionnel}

\begin{itemize}
\tightlist
\item
  \texttt{modele}, \texttt{equation}, \texttt{r2},
  \texttt{asymptote\_y}, \texttt{asymptote\_x}
\end{itemize}

\paragraph{\texorpdfstring{\texttt{ICR\_06\_occupation\_spatiale.csv}
(conditionnel)}{ICR\_06\_occupation\_spatiale.csv (conditionnel)}}\label{icr_06_occupation_spatiale.csv-conditionnel}

\begin{itemize}
\tightlist
\item
  \texttt{espece}, \texttt{nb\_placettes}, \texttt{prop\_placettes}
\end{itemize}

\paragraph{\texorpdfstring{\texttt{ICR\_07\_metrics\_pertinence.csv}}{ICR\_07\_metrics\_pertinence.csv}}\label{icr_07_metrics_pertinence.csv}

\begin{itemize}
\tightlist
\item
  \texttt{site}, \texttt{nb\_visites}, \texttt{terminal\_window\_k},
  \texttt{tee\_final}, \texttt{ir\_final}, \texttt{completude},
  \texttt{r2\_lineaire}, \texttt{r2\_hyperbolique},
  \texttt{slope\_terminal\_cumul}, \texttt{tail\_novelty\_share},
  \texttt{spearman\_temporal\_spatial}, \texttt{discordance\_rate},
  \texttt{score\_pertinence}, \texttt{class\_ir},
  \texttt{class\_completude}
\end{itemize}

\begin{center}\rule{0.5\linewidth}{0.5pt}\end{center}

\subsection{11. Métriques et
interprétation}\label{muxe9triques-et-interpruxe9tation}

\subsubsection{10.1 Seuils de classes
implémentés}\label{seuils-de-classes-impluxe9mentuxe9s}

\texttt{class\_ir} :

\begin{itemize}
\tightlist
\item
  \texttt{\textless{}\ 0.60} : \texttt{faible}
\item
  \texttt{{[}0.60,\ 0.80)} : \texttt{moyenne}
\item
  \texttt{\textgreater{}=\ 0.80} : \texttt{bonne}
\end{itemize}

\texttt{class\_completude} :

\begin{itemize}
\tightlist
\item
  \texttt{\textless{}\ 0.70} : \texttt{insuffisante}
\item
  \texttt{{[}0.70,\ 0.90)} : \texttt{intermédiaire}
\item
  \texttt{\textgreater{}=\ 0.90} : \texttt{avancée}
\end{itemize}

\subsubsection{10.2 Lecture recommandée des
indicateurs}\label{lecture-recommanduxe9e-des-indicateurs}

\begin{itemize}
\tightlist
\item
  \texttt{Ir\ final} élevé + \texttt{tail\_novelty\_share} faible :
  inventaire bien stabilisé.
\item
  \texttt{completude\_obs\_sur\_asymptote} proche de 1 : effort proche
  de l'asymptote estimée.
\item
  \texttt{spearman\_temporal\_spatial} élevé et
  \texttt{discordance\_rate} faible : cohérence temporelle/spatiale
  satisfaisante.
\end{itemize}

\begin{center}\rule{0.5\linewidth}{0.5pt}\end{center}

\subsection{12. Versioning et
maintenance}\label{versioning-et-maintenance}

\subsubsection{11.1 Version du script}\label{version-du-script}

La version est portée dans l'en-tête du fichier script
(\texttt{Version\ :\ 1.4}).

\subsubsection{11.2 Maintenance
recommandée}\label{maintenance-recommanduxe9e}

\begin{itemize}
\tightlist
\item
  maintenir la cohérence entre \texttt{README.md} et ce document ;
\item
  mettre à jour cette doc si les noms de sorties \texttt{ICR\_*}
  évoluent ;
\item
  conserver les seuils de classes synchronisés avec
  \texttt{calc\_scientific\_metrics()}.
\end{itemize}

\begin{center}\rule{0.5\linewidth}{0.5pt}\end{center}

\subsection{13. Exécution, dépendances et
limites}\label{exuxe9cution-duxe9pendances-et-limites}

\subsubsection{12.1 Dépendances}\label{duxe9pendances}

\textbf{Requises} :

\begin{itemize}
\tightlist
\item
  \texttt{dplyr}, \texttt{tidyr}, \texttt{ggplot2}, \texttt{purrr},
  \texttt{readr}, \texttt{stringr}, \texttt{forcats}, \texttt{tibble},
  \texttt{scales}, \texttt{gridExtra}
\end{itemize}

\textbf{Optionnelles} :

\begin{itemize}
\tightlist
\item
  \texttt{minpack.lm} (modèle hyperbolique)
\item
  \texttt{vegan} (CA/AFC)
\end{itemize}

\subsubsection{12.2 Exécution}\label{exuxe9cution}

\begin{verbatim}
Rscript scripts/Inventaires_completude_representativite.R
\end{verbatim}

Désactiver l'auto-exécution :

\begin{verbatim}
INVENTAIRES_AUTO_RUN=FALSE Rscript scripts/Inventaires_completude_representativite.R
\end{verbatim}

\subsubsection{12.3 Limites connues}\label{limites-connues}

\begin{itemize}
\tightlist
\item
  sans \texttt{placette} : pas d'analyse spatiale dédiée ;
\item
  sans \texttt{minpack.lm} : complétude asymptotique potentiellement
  \texttt{NA} ;
\item
  dates non conformes à \texttt{CONFIG\$date\_format} : arrêt de
  l'exécution ;
\item
  en mode strict, un CSV mal formé est bloquant (par design qualité) ;
\item
  faible nombre de visites : qualité des modèles limitée.
\end{itemize}

\begin{center}\rule{0.5\linewidth}{0.5pt}\end{center}

\subsection{14. Annexes}\label{annexes}

\subsubsection{13.1 Glossaire}\label{glossaire}

{\def\LTcaptype{none} % do not increment counter
\begin{longtable}[]{@{}
  >{\raggedright\arraybackslash}p{(\linewidth - 2\tabcolsep) * \real{0.5000}}
  >{\raggedright\arraybackslash}p{(\linewidth - 2\tabcolsep) * \real{0.5000}}@{}}
\toprule\noalign{}
\begin{minipage}[b]{\linewidth}\raggedright
Terme
\end{minipage} & \begin{minipage}[b]{\linewidth}\raggedright
Définition
\end{minipage} \\
\midrule\noalign{}
\endhead
\bottomrule\noalign{}
\endlastfoot
Visite distincte & Couple \texttt{date\ +\ visite\_id} (et \texttt{site}
en multi-sites) \\
TEE & Taux d'espèces observées une seule fois \\
Ir & Indice de représentativité (\texttt{1\ -\ TEE}) \\
Complétude & Ratio richesse observée / asymptote hyperbolique \\
Occupation spatiale & Proportion de placettes où l'espèce est
observée \\
Score de pertinence & Score synthétique pondéré multi-indicateurs \\
vegan & Package R d'analyse en écologie des communautés (ordinations
CA/AFC, ACP, NMDS ; indices de diversité ; tests de permutation).
Utilisé ici pour l'AFC placettes-espèces si \texttt{make\_ca\ =\ TRUE}
et colonne \texttt{placette} disponible. \\
\end{longtable}
}

\subsubsection{13.2 Résumé ultra-compact (1
tableau)}\label{ruxe9sumuxe9-ultra-compact-1-tableau}

{\def\LTcaptype{none} % do not increment counter
\begin{longtable}[]{@{}
  >{\raggedright\arraybackslash}p{(\linewidth - 6\tabcolsep) * \real{0.2500}}
  >{\raggedright\arraybackslash}p{(\linewidth - 6\tabcolsep) * \real{0.2500}}
  >{\raggedright\arraybackslash}p{(\linewidth - 6\tabcolsep) * \real{0.2500}}
  >{\raggedright\arraybackslash}p{(\linewidth - 6\tabcolsep) * \real{0.2500}}@{}}
\toprule\noalign{}
\begin{minipage}[b]{\linewidth}\raggedright
Type
\end{minipage} & \begin{minipage}[b]{\linewidth}\raggedright
Fichier
\end{minipage} & \begin{minipage}[b]{\linewidth}\raggedright
Condition
\end{minipage} & \begin{minipage}[b]{\linewidth}\raggedright
Finalité
\end{minipage} \\
\midrule\noalign{}
\endhead
\bottomrule\noalign{}
\endlastfoot
Entrée & \texttt{data/observations.csv} & Obligatoire & Source des
observations \\
CSV global & \texttt{ICR\_donnees\_preparees.csv} & Toujours & Base
nettoyée \\
CSV global & \texttt{ICR\_resume\_tous\_sites.csv} & Toujours & Synthèse
par site \\
CSV global & \texttt{ICR\_metrics\_pertinence\_tous\_sites.csv} &
Toujours & Métriques consolidées \\
CSV site & \texttt{ICR\_01} à \texttt{ICR\_07} (\texttt{*.csv}) & Selon
cas & Détails analytiques par site \\
PNG global & \texttt{ICR\_comparaison\_*.png} & Toujours & Comparaisons
inter-sites \\
PNG site & \texttt{ICR\_01} à \texttt{ICR\_07} (\texttt{*.png}) & Selon
cas & Visualisations locales \\
\end{longtable}
}
